\documentclass[12pt]{article}
\usepackage[utf8]{inputenc}
\usepackage[T1]{fontenc}
\usepackage[margin=1in]{geometry}
\usepackage{setspace}
\usepackage{xcolor}
\usepackage[colorlinks=true, urlcolor=blue!60!black, linkcolor=blue!60!black]{hyperref}
\setlength{\parindent}{0em}
\setlength{\parskip}{0.6em}
\renewcommand{\baselinestretch}{1.0}
\begin{document}
\begin{center}
{\LARGE \textit{Heart of Atlanta Motel, Inc.\ v.\ United States}, 379 U.S.\ 241 (1964)\par}
\vspace{0.5em}
{\large Case Brief --- Two-Page Edition\par}
\end{center}

\section{Memory Jogger}
Congress may use the Commerce Clause to ban racial discrimination by motels serving interstate travelers, even though its aim is moral (\href{https://supreme.justia.com/cases/federal/us/379/241/}{Justia}).

\section{Detailed Case Facts}
Title II of the Civil Rights Act of 1964 barred discrimination in public accommodations affecting commerce, including motels serving transient guests. The Heart of Atlanta Motel, near two interstates, advertised nationally, drew about 75\% of its guests from out of state, and refused Black guests.

\section{Procedural History}
A three-judge district court upheld Title II and enjoined the motel; the Supreme Court affirmed on direct appeal.

\section{Judicial Votes}
9--0. Clark wrote for eight Justices; Black concurred in the judgment and holding; Douglas and Goldberg joined and also concurred.

\section{Holding}
Congress may prohibit discrimination by public accommodations serving interstate travelers when it rationally finds an effect on commerce and uses reasonable means, even while addressing a moral wrong. Title II does not violate the Fifth or Thirteenth Amendment.

\section{Analysis of Opinions}
\textbf{Majority:} Unlike the 1875 Act struck in the \textit{Civil Rights Cases}, Title II targets enterprises tied to interstate commerce. The hearings showed that discrimination deterred Black travel; the movement of persons is commerce, and Congress may reach its local obstructions. \textbf{Black:} aggregation under \textit{Wickard} sustains Title II, though remote effects would not. \textbf{Douglas and Goldberg:} the Act should also rest on \S~5 of the Fourteenth Amendment, since its purpose is human dignity.

\section{Examples: Future Applications}
\textbf{Same-Side:} barring discrimination by commercial vacation-rental hosts serving travelers. \textbf{Opposite-Side:} a remedy for local gender-motivated violence fails under \textit{Morrison}. \textbf{Fence-Sitter:} extending Title II to a remote two-bedroom farmhouse with only in-state guests.

\section{Critique}
Originalists and libertarians such as Randy Barnett and Richard Epstein say the decision made the Commerce Clause a police power; progressives regret that it framed civil rights in commercial terms.

\section{Key Quotations}
\begin{quote}
``Thus, the power of Congress to promote interstate commerce also includes the power to regulate the local incidents thereof, including local activities in both the States of origin and destination, which might have a substantial and harmful effect upon that commerce.''
\end{quote}
\noindent --- Justice Clark, for the Court, 379 U.S.\ at 258.

\end{document}
